\section{Team organization and development process}

\subsection{Division of work}
\begin{table}[htbp]
\centering
\footnotesize
\begin{tabularx}{\textwidth}{@{}p{18mm}p{34mm}p{34mm}X@{}}
\toprule
\textbf{Student ID} & \textbf{Member} & \textbf{Primary module} &
\textbf{Principal contribution evidence} \\
\midrule
24125023 & Nguyen Trong Van Viet & Destinations and quests & Versioned destination
schema; Nha Trang and Ho Chi Minh City packs; Explore, route detail, Free Run,
proximity/fallback rules, and related tests \\
24125078 & Nguyen Hong Tan Tai & Tracking and optional sharing & Run lifecycle,
location stream, distance and jump filtering, duration and pace, checkpoints,
foreground settings, and live-share adapter \\
24125093 & Nguyen Dinh Thien Loc & Team leader; Story Studio, Tracks, and report &
Product coordination, integration and release review; adaptation of the
MIT-licensed Lockit camera baseline into the Tracks navigation, offline capture,
mock friend and sharing flow; report organization and production; serializable
story model, original templates, element operations, undo/redo, exact PNG
renderer, Android sharing, and related tests \\
24125107 & Tran Le Anh Tuan & Persistence, History, integration & SQLite schema and
repositories, retained media, deletion cascade, achievements, Journey states,
dependency assembly, recovery, Android configuration, and persistence tests \\
\bottomrule
\end{tabularx}
\caption{Team ownership and observable deliverables.}
\end{table}


Nguyen Dinh Thien Loc served as the team leader. In addition to owning Story
Studio and export, Loc adapted and integrated the inherited Lockit camera into
the production Tracks flow, implemented its offline-safe capture and mock social
sharing states, coordinated the shared product direction, aligned work at
integration points, reviewed release readiness, and organized and produced the
final report and demonstration build. Leadership was collaborative: module
owners retained responsibility for their deliverables, while cross-module
decisions were made against the same offline-first architecture and acceptance
criteria.

The repository history contains commits attributed to all four members. Module
ownership identifies responsibility rather than exclusive access: each member
could review other modules, while changes to domain contracts, dependency
configuration, manifests, navigation, and the package specification were treated
as integration hotspots.

\subsection{Collaboration method}
Work was divided against interfaces so a feature did not need to wait for a
database, device, or another unfinished screen. Every principal contract has a
deterministic fake. This made it possible to develop and test routes, tracking,
Story Studio, History, and optional sharing independently before integration.

The shared engineering contract required contributors to inspect existing work,
preserve unrelated changes, use stable domain language, keep remote behavior
optional, format changed code, run proportionate tests, and state verification
limits. Suggested feature branches and focused commits preserved visible member
history. Secrets, private tracks, build caches, database files, and generated
indexes were excluded from source control.

\subsection{Integration decisions}
The team selected immutable shared models before creating feature variants.
Presentation depends on contracts, while SQLite and device plugins remain data
adapters. Destination content is supplied through a repository contract;
tracking persists through the run contract; Story Studio consumes a completed
run summary; and History combines the repositories through dependency assembly.

This approach reduced merge coupling and made failure behavior testable. It also
made architectural debt visible: cross-feature navigation and shared visual
constants are integration responsibilities, not reasons for one feature to reach
into another feature's private implementation.
